\begin{center}
\LARGE
About an Algorithm to Construct a Canonical Basis for the
           Ideals in the Ring of Formal Power Series

\end{center}

\begin{center}
\Large
Nikolay N. Vasiliev, Victor Edneral

\end{center}

\noindent
Date:  July 19th (Friday)\\
Time:  15:10-15:30 \\

\begin{center}
\large
Abstract
\end{center}

It is well known fact that the problem to solve systems of polynomial algebraic equations can be
reduced to the problem of construction of Groebner basisis for the ideals in the ring of
polynomials. In the case of the equations in the ring of formal power series this technic does not
work because one of the important thing for Buchberger's algorithm is a definition of "leading
term" of polynomial. It does not exist in the case of formal power series. 

>From other side the problem to solve systems of equations for the finite parts of the infinite
formal series is very important for different applications. For example we could say about
investigations of singular points of the systems of ODE. Another example is the problem of
normalization of the ODE. [1,2]. 

Usually one can consider the solutions for the finite parts of the series only. It allow us to consider
the system of the equations over the formal power series ring as usuall polynomial system, but this
approach does not lead to good result, because adding terms of highest order could change your
solutions in lower terms. 

We can show that some modification of Buchberger's algotithm allow us to construct a canonical
basis whis is full analogue of the GB in the case of ideas in the ring of formal power series. The
construction of such canonical basis and the proof of appropriate variant of the Buchberger's
theorem about finiteness of the algorithm is very important not only in the connection with the
problem to solve the equations in the ring of formal power series but for studing of the structure of
the ring of the FPS and it's ideals. This algorithm get us a canonical representation for the ideals of
this ring. 

We use in our algorithm some special representation for basic elements of the ideal. This
representation uses monomials which lie on the border of specially modified Newton polytopes for
the series under consideration. For each serie of the system of we have only a finite number of such
monomials. If we have the system of finite number of "basic" monomials we can represent the
serie as a sum such monomials with coefficients which are also the series. It is very important
thing for the algorithm that these coefficients are invertible elements of the ring of FPS. As a
result we can evaluate by usual method S polynomials and make reductions of the leading terms
only for Newton polytops. It means that each modified Newton polytop of the series of the
equations is irreducible relatively others modified Newton Polytops. We suggest to use such sets
of FPS as standart basisis in the ring of FPS. Such sets of FPS have many qualities of usual
Groebner basisis in the case of polynomials. The proofs of the variants of the Buchberger's
theorems are based on the fact that the ring of FPS over a noetherian commutative ring is
noetherian also. 

\bigskip
\noindent 
{\bf References}
\begin{enumerate}

\item
 Bruno A. D. ``Local Method in Nonlinear Differential Equations.'' Part I-The Local Method of
Nonlinear Analyses of Differential Equations, Part II-The Sets of Analyticity of a Normalizing
Transformation. Springer Series in Soviet Mathematics, 1989. ISBN 3-540-18926-2, 370 p. 


\item
Edneral V. F. ``Computer Generation of Normalizing Transformation for Systems of Nonlinear
ODE.'' Proceedings of the 1993 International Symposium on Symbolic and Algebraic
Computation (Kiev, Ukraine, July 6-8, 1993) / M. Bronstein. N.Y.: ACM Press, 1993. P. 14--19. 


\item
 N.N.Vassiliev, V.F.Edneral, ``About an Algorithm to Construct a Canonical Basis for the Ideals
in the Ring of Formal Power Series'', in the book of Abstracts of the Russian Conference
"Computer Methods of Celestial Mechanics -95", St. Petersburg, 1995, Ed. by prof Sokolsky. 


\end{enumerate}